\section{Experiments}
\label{sec:experiments}

\subsection{Implementation Details}

\noindent\textbf{Base Generation \& Cascade Configuration.}
To maximize the diversity and quality of the initial candidate pool ($C_{init}$), our base generation utilizes a strict Multi-Provider Cascade. We query foundation models in increasing order of spatial reasoning capability: OpenAI Codex (\texttt{gpt-5.3-codex}), followed by OpenAI's core model (\texttt{gpt-5.3})~\cite{openai2025codex}, and finally handled by Anthropic (\texttt{claude-4.6-Opus}~\cite{anthropic2026opus}) for the most complex prompts. The input prompts include the full Ops JSON string—which is also generated by the VLM for human-readable semantic decomposition—alongside rigorous instructions restricting the output strictly to valid CadQuery syntax. To encourage deterministic geometric generation over creative hallucination, we set the sampling temperature to a strictly low threshold ($T=0.1$).

\noindent\textbf{Evaluation Infrastructure.}
All geometric compilations, volumetric measurements, and exact Boolean intersections are executed utilizing the OpenCASCADE (OCCT) Python kernel. The 3D visual feedback provided to the Re-Act baselines consists of orthogonal 4-view renderings (Front, Top, Right, Isometric) generated automatically from the compiled STEP artifacts.

\subsection{Data Distillation Yield \& Proximal Optimization}

We evaluated the overarching pipeline on a large corpus of 1,619 parts from the Fusion360 Gallery. To explicitly isolate the effectiveness of our parametric grounding, we focus our primary analysis on the proximal boundary subset ($N=226$), where the Direct VLM generation achieves an initial topological accuracy of $\text{IoU} \ge 0.8$.

\begin{table}[h]
  \centering
  \small
  \setlength{\tabcolsep}{4pt}
  \caption{\textbf{End-to-End Generative Yield.} Evaluated on candidates with initial viable topologies ($N=226$). CADLoop achieves absolute convergence while driving geometric error to near-zero.}
  \label{tab:pipeline}
  \begin{tabular}{@{}lccc@{}}
    \toprule
    \textbf{Method} & \textbf{Resolved} & \textbf{Pass Rate} ($\uparrow$) & \textbf{CD (mm)} ($\downarrow$) \\
    \midrule
    Direct VLM      & 194 & 85.8\% & 7.07 \\
    + Re-Act        & 207 & 91.6\% & 3.36 \\
    \textbf{+ CADLoop} & \textbf{226} & \textbf{100\%} & \textbf{$<$0.01} \\
    \bottomrule
  \end{tabular}
\end{table}

\noindent\textbf{Defending Precision over Recall.}
While a 14.6\% global one-shot pass rate (across the entire 1,619 corpus) strongly reflects the extreme zero-shot difficulty of unconstrained CAD generation, it embodies a critical ``Precision over Recall'' curation philosophy. Forcing a probabilistic pipeline to hallucinate ``fixes'' for structurally collapsed cases ($\text{IoU} < 0.5$) inevitably pollutes the dataset with structurally unsound, nonsensical scripts. By aggressively discarding irrecoverable topologies and focusing solely on high-confidence parametric refinements, CADLoop distills a pristine, noise-free dataset. Crucially, as shown in \cref{tab:pipeline}, on the accepted viable files, CADLoop successfully reduces the final mean geometric error (Chamfer Distance) by over 99\%, achieving near-perfect reconstruction ($<0.01$ mm).

\subsection{Ablation Study: Isolating the Fallback Mechanism}

To rigorously prove that the Unified Geometric Skill out-performs state-of-the-art probabilistic self-correction, our ablation explicitly isolates the performance of Direct VLM, Re-Act, and CADLoop within specific geometric failure regimes. Crucially, out of the 226 proximal boundary candidates, 194 passed initially, leaving exactly 32 critical near-miss failures for the system to repair.

\begin{table}[h]
  \centering
  \small
  \setlength{\tabcolsep}{4pt}
  \caption{\textbf{Ablation of Repair Efficacy on Failure Cases.} Evaluated on exact failures per regime. Re-Act hits a reasoning ceiling on near-misses; CADLoop eliminates them entirely.}
  \label{tab:ablation}
  \begin{tabular}{@{}lcccc@{}}
    \toprule
    \textbf{B0 IoU Regime} & $N$ & \textbf{Re-Act} & \textbf{CADLoop} & \textbf{$\Delta$} \\
    \midrule
    Structural Plateau (0.50--0.80) & 38     & 0.0\%  & 2.6\%  & +2.6\% \\
    Near-miss (0.80--0.99)          & 32$^*$ & 40.6\% & \textbf{100.0\%} & \textbf{+59.4\%} \\
    \bottomrule
    \multicolumn{5}{l}{\scriptsize $^*$ 32 near-miss failures from the $N{=}226$ viable set ($\text{B0} \ge 0.8$); 194 already pass.}
  \end{tabular}
\end{table}

\noindent\textbf{The Re-Act Bottleneck.}
As demonstrated in \cref{tab:ablation}, despite having access to high-fidelity multi-view renders and exact IoU numerical feedback, the Re-Act baseline hits a severe reasoning ceiling. It resolves only 13 out of the 32 near-misses ($40.6\%$). Visual feedback proves insufficient if the neural agent fundamentally lacks the internal capacity for continuous mathematical grounding. The agent frequently guesses random radius or offset values, failing to converge.

\noindent\textbf{The Dominance of the Unified Skill.}
By replacing the probabilistic Re-Act loop with CADLoop's deterministic fallback skill and manual flywheel, the system achieves a remarkable 100\% repair rate on the 32-part proximal boundary subset (a massive \textbf{+59.4\%} absolute gain over Re-Act). This validates our core hypothesis: once a VLM establishes a structurally sound sketch skeleton, explicit programmatic tools—not more neural guessing—are the optimal mechanism for securing spatial equivariance.

\subsection{Qualitative Analysis}

\begin{figure*}[t]
  \centering
  \includegraphics[width=\textwidth]{figure_4}
  \caption{%
    \textbf{Auto-Pass Examples via CADLoop.} 
    The fallback skill successfully resolves spatial mapping errors on complex, real-world geometries, including intricate gear teeth arrays (left) and extruded connector flanges (right), completely without multi-modal VLM re-prompting overhead.
  }
  \label{fig:autopasses}
\end{figure*}

As illustrated in \cref{fig:autopasses}, CADLoop robustly handles immense geometric diversity. The leftmost object features complex rotational symmetries (gear teeth), while the subsequent objects feature precise angular cutouts and multi-plane extrusion channels. CADLoop's automated coordinate-plane matching effortlessly resolves initial axis-inversion errors.

\begin{figure*}[t]
  \centering
  \includegraphics[width=\textwidth]{figure_5}
  \caption{%
    \textbf{Hard Case Resolution Analysis.} 
    Red dotted boxes denote specific geometric discrepancies. Middle column: the VLM inverted a fillet extrusion direction, ignoring semantic text clues. Rightmost column: the VLM hallucinated a filled solid block instead of properly parsing the \texttt{is\_outer=False} logic for an inner hole cutout. CADLoop deterministically patches the AST to explicitly fix these boundary mismatches.
  }
  \label{fig:hardcases}
\end{figure*}

\Cref{fig:hardcases} showcases the resolution of persistent hard cases that completely break standard VLMs. In the ring structure (middle), the base model generated a concave sweep instead of a convex fillet. While Re-Act struggles to distinguish $+Z$ from $-Z$ based on visual feedback, CADLoop mathematically intersects the expected bounding boxes to force the correct extrusion sign. In the angled bracket (right), a critical functional hole was omitted. The CADLoop AST-patching algorithm isolates the missing profile loop and mathematically splices a precise Boolean cut command into the final script, converting a failed generation into a perfect SFT asset.